\scanpage{132}
\mspage{77}
% archive sheet number: 131

établissent des équivalences quasi-inverses l'une de l'autre entre

\diagram{\begin{tikzcd}[ampersand replacement=\&,column sep=large]
{\unsure{\mathcal{F}|F}'\Sigma^{-1}}\arrow[r,"\approx"]\&\mathcal{F}\Sigma^{-1}\arrow[l,bend left=15,"{\underset{\approx}{\Phi_{\mathcal{F}',\mathcal{F}}}}"]\&(\Sigma=W_S\text{ ou }W_S^d).
\end{tikzcd}}

Énoncé dual pour le foncteur

\[
\Psi:\mathrm{Cat}_{/S}\longrightarrow\underline{\mathrm{Cofibsc}}\,S,
\]

associant à $X$ de $\mathrm{Cat}_{/S}$ la catégorie cofibrée scindée sur $S$

\[
s\longmapsto X_{/s}.
\]

Il donne des \unsure{couples} de foncteurs quasi-inverses

\diagram{\begin{tikzcd}[ampersand replacement=\&,column sep=large]
(\underline{\mathrm{Cofibsc}}\,S)\Sigma^{-1}\arrow[r]\&\mathcal{F}\Sigma^{-1}\arrow[l,bend left=14,"\Psi_{\mathcal{F}}"]\&(\Sigma=W_S\text{ ou }W_S^g),
\end{tikzcd}}

pour toute sous-catégorie pleine \inserted{$\mathcal{F}$} de $\mathrm{Cat}_{/S}$ telle que

\[
\underline{\mathrm{Cofib}}_S\subset\mathcal{F}\subset\mathrm{Cat}_{/S},
\]

et des couples de foncteurs quasi-inverses

\diagram{\begin{tikzcd}[ampersand replacement=\&,column sep=large]
\mathcal{F}'\Sigma^{-1}\arrow[r,"\approx"]\&\mathcal{F}\Sigma^{-1}\arrow[l,bend left=15,"\approx"',"{\Psi_{\mathcal{F}',\mathcal{F}}}"]\&(\Sigma=W_S\text{ ou }W_S^g)
\end{tikzcd}}

pour deux sous-catégories pleines $\mathcal{F}',\mathcal{F}$

\[
\underline{\mathrm{Cofib}}_S\subset\mathcal{F}'\subset\mathcal{F}\subset\mathrm{Cat}_{/S}.
\]

Les foncteurs $\Phi,\Psi$ \deleted{induisent sur} \inserted{envoient} la sous-catégorie pleine $\underline{\mathrm{Fibhot}}\,S$ \deleted{\unsure{des}} ``l'étage'', dans les sous-catégories $\underline{\mathrm{Fibsc}}_0\,S$, $\underline{\mathrm{Cofibsc}}_0\,S$ respectivement, donc induisent à l'étage des foncteurs $\Phi_0,\Psi_0$

\diagram{\begin{tikzcd}[ampersand replacement=\&,column sep=large]
\deleted{\unsure{\Phi_0}}\underline{\mathrm{Fibsc}}_0\,S\&\underline{\mathrm{Fibhot}}\,S\arrow[l,"\Phi_0"']\arrow[r,"\Psi_0"]\&\underline{\mathrm{Cofibsc}}_0\,S,
\end{tikzcd}}

jouant à l'étage exactement le $=$ rôle que $\Phi,\Psi$ \unsure{au|en} rez de chaussée. La seule différence, c'est qu'à l'étage on a $W_S=W_S^d=W_S^g\deleted{\unsure{=W_S}}\,(=W_S^{\unsure{f|t}})$,
